\documentclass[11pt]{article}
\usepackage{amsmath,amssymb,amsthm}
\newtheorem{theorem}{Theorem}
\newtheorem{lemma}{Lemma}
\begin{document}
\title{Independent Audit of Reverse-Orientation Gluing}
\author{Research run: Reverse Orientation Head Compression}
\maketitle

\paragraph{Status.}
This note audits the reduction in \texttt{exponent-mechanisms.tex}
and \texttt{exponent-literature.tex}. Its statements are unconditional
finite implications with explicit hypotheses. They do not establish the
missing universal local compiler or the requested new asymptotic bound.
The unchanged model admits every affine ratio with positive constant
coefficient and strictly positive slopes, by
\texttt{FracAtom.ofPositiveAffineRatio\_eval}, and every everywhere-positive
affine ratio with strictly negative slopes, by
\texttt{FracAtom.ofNegativeAffineRatio\_eval}. Different heads may have
different orientations. The proofs below use these existing interfaces.

\section{Gluing}
Identify a Boolean input with $(X,Y)$ in two disjoint coordinate blocks
$A,B$. For a center $(S,T)$ define
\[
D_{S,T}=\{(X,Y):X\subseteq S,\ Y\subseteq T\},\qquad
C_{S,T}=\{(S\setminus\{i\},T\setminus\{j\}):i\in S,\ j\in T\}.
\]
The centers form a finite list; repetitions are harmless. Empty cells
are discarded. All labels are restrictions of one fixed function
$f:\{0,1\}^n\to\{-1,1\}$, so overlapping cells have consistent labels.

\begin{theorem}
Suppose each center has a bias plus $H_{S,T}$ affine ratios with
positive constant coefficients and nonnegative slopes in the original
$x$ coordinates. Suppose its score $F_{S,T}$ has the required strict
signs on $C_{S,T}$ and vanishes on all points of $D_{S,T}$ where
$X=S$ or $Y=T$. Then a score with at most
$\sum H_{S,T}$ admissible heads has those strict signs on the union
of the cells.
\end{theorem}
\begin{proof}
A nonempty cell and the axis-zero hypothesis force $H_{S,T}\geq1$:
with no ratios the score is constant, zero on the axes, and cannot have
a strict sign anywhere. Absorb its bias $c$ into the first numerator
as $A_1\leftarrow A_1+cB_1$. This changes neither the score nor the
number of heads. Crucially this is done before the following penalty.

Process centers in nonincreasing order of $|S|+|T|$. Given the current
score $G$ and its protected set $P$, write
\[
q(x)=\sum_{i\in A\setminus S}x_i+\sum_{j\in B\setminus T}x_j,
\qquad
F_t(x)=\sum_h\frac{A_h(x)}{B_h(x)+tq(x)},\quad t\geq0.
\]
Every denominator has its original positive constant coefficient and
nonnegative slopes. On $D_{S,T}$ one has $F_t=F_{S,T}$ exactly.
At every point outside $D_{S,T}$, $q(x)\geq1$, so $F_t(x)\to0$.

Choose $\lambda>0$ so that
$f(x)(G(x)+\lambda F_{S,T}(x))>0$ on the new cell. This is possible
because the cell is finite and all the coefficients
$f(x)F_{S,T}(x)$ are strictly positive.
An earlier protected point inside $D_{S,T}$ belongs to a preceding
cell with center rank at least $|S|+|T|$. Its vertex rank is at least
$|S|+|T|-2$. Thus its two nonnegative deletion counts sum to at most
two. If one count is zero, the new score is zero there. If both are
positive, both are one, and the point lies in the new cell. Its new
contribution has the same strict sign as its protected score. This
argument includes ties in center rank and overlapping cells.

At the remaining protected points the finite positive margin of $G$
allows one common sufficiently large $t$ so that
$|\lambda F_t(x)|<|G(x)|$ at every such point. Hence all earlier
strict signs persist. Scaling each numerator by $\lambda$ uses
exactly $H_{S,T}$ new ratios. Induction completes the weak-slope sum.

Finally replace every denominator in the completed sum by
$B(x)+\delta\sum_i x_i$ with $\delta>0$. Every slope is now strictly
positive. Every denominator retains its positive constant coefficient.
The finite sum converges pointwise, hence uniformly on the finite cube,
to the unperturbed sum as $\delta\to0$. Choose a common $\delta$ that
preserves the positive minimum margin on the covered set. This gives
actual admissible atoms with no head-count increase. Axis zeros may
cease to hold at this last step; they are no longer needed.
\end{proof}

\paragraph{Local-to-ambient extension.}
For a local denominator $b_0-\sum_i p_i z_i-\sum_jq_jw_j$ with
$p_i,q_j\geq0$, positivity on the full local cube implies
$b_0-\sum_ip_i-\sum_jq_j>0$. Substitute $z_i=1-x_i$ and $w_j=1-x_j$.
The resulting constant coefficient is precisely this positive number,
and all its slopes are nonnegative. Add any nonnegative slopes on
coordinates outside $S\cup T$. Those coordinates vanish throughout
$D_{S,T}$, so the local scores and axis zeros are unchanged there.
All resulting denominators are positive on the whole ambient cube.
Thus positivity is checked globally, not merely on the promised cell.

\section{Cover and head count}
The existing declarations \texttt{monomialStarCoverNumber\_spec} and
\texttt{monomialStarCoverNumber\_real\_le\_eleven} supply, for each
$m\geq1$, an upward closed-star cover of size at most $11\cdot2^m/m$.
Complementing every set gives a downward closed-star cover of the same
size. For each of its centers $S\ne[m]$, choose one $i_S\notin S$
and add $S\cup\{i_S\}$. Keep all original centers. This increases
the size by at most a factor two. Every proper vertex is now in an
open downward star: if it was already a child, use its old center;
if it was only the center itself, use the chosen larger center.
The full set cannot be an open downward child of any center.

For $a=\lfloor n/2\rfloor$, $b=\lceil n/2\rceil$, $n\geq2$, the
products of these open covers therefore cover exactly
\[
U=(\{0,1\}^a\setminus\{\mathbf1\})
\times(\{0,1\}^b\setminus\{\mathbf1\})
\]
with at most
\[
N\leq\frac{484\cdot2^n}{ab}=O(2^n/n^2)
\]
nonempty modules. Repeated centers can be removed. If a local compiler
uses at most $K(1+m/\log\log(m+4))$ heads, define
\[
M_n=K\max_{2\leq m\leq n}\left(1+\frac{m}{\log\log(m+4)}\right).
\]
The gluing budget is at most $NM_n$. For $m\leq\sqrt n$, positivity
of $\log\log6$ gives $O(\sqrt n)$; for $m>\sqrt n$, use
$\log\log(m+4)\geq\log\log n-\log2$. Hence
$M_n=O(n/\log\log n)$. The same split gives $O(n/\log n)$ for a
local $O(1+m/\log(m+4))$ compiler.

\section{Exceptional faces}
Let $G$ be the actual admissible glued score on $U$. For the first
exceptional face $X=A$, interpolate the real function
$Y\mapsto f(A,Y)-G(A,Y)$ by its unique multilinear polynomial $p(Y)$.
It has at most $2^b$ monomials. Add the formal score
$x_Ap(Y)$, where $x_A=\prod_{i\in A}x_i$. This changes only the first
face and makes its values exactly $f$. On the second face $Y=B$,
interpolate the residual error in $X$ by a polynomial $q(X)$, using
at most $2^a$ monomials, and add $x_Bq(X)$. Its error at $X=A$ is
zero because the intersection has already been corrected. The second
correction is therefore zero on the entire first face: away from
$Y=B$ its multiplier is zero, and at the intersection $q(A)=0$.
The resulting target score has correct strict signs on the whole cube.

For every coordinate set $V$, the single admissible ratio
\[
Q_{V,\eta}(x)=\frac{\eta}{\eta+\sum_{i\in V}(1-x_i)
+\eta^2\sum_{i\notin V}(1-x_i)},\qquad\eta>0,
\]
has strictly negative slopes and denominator at least $\eta$.
Where $x_V=0$ it is at most $\eta$; where $x_V=1$ its difference
from one is at most $n\eta$. Thus for $n\geq1$ the uniform error
is at most $n\eta$. This includes $V=\varnothing$.
If the correction coefficients are $d_V$, the total approximation
error is at most $n\eta\sum_V|d_V|$. Choose $\eta$ small relative
to the positive minimum strict margin of the formal target score.
Absorb each $d_V$ in the constant numerator of its ratio. At most
$2^a+2^b\leq2\cdot2^{\lceil n/2\rceil}$ extra admissible heads
therefore suffice. No new input coordinates are introduced.

The full conditional budget, with every term accounted for, is
\[
H\leq\frac{484\cdot2^n}{ab}\,M_n+2^a+2^b.
\]
The last term is $o(2^n/(n\log\log n))$, and is also
$o(2^n/(n\log n))$. Bias absorption, suppression, extension, and the
final strict-slope perturbation add no heads. The local compiler is
still an external unproved hypothesis; none of this audit supplies it.
\end{document}
